
Reinforcement Learning from Human Feedback (RLHF) (Ouyang et al., 2022; Bai et al., 2022) is a widely used method for aligning language model outputs with human preferences. Standard RLHF evaluation treats the interaction as unidirectional: human preference annotations shape AI outputs, and evaluation metrics focus on AI response quality. Whether and how AI response quality is associated with downstream human communicative behavior in these conversations has received comparatively little empirical attention.

Communication Accommodation Theory (CAT) \citep{Giles1973Accent} predicts that interlocutors adapt their communication patterns toward one another, with lower-power interlocutors tending to accommodate more to higher-power ones \citep{DanescuNiculescuMizil2012Echoes}. Prior empirical work on accommodation has primarily operated at the lexical level (function-word coordination) in human-human settings. Recent work has examined word-level style adaptation in human-AI conversations \citep{Chang2025Language} and proposed conceptual frameworks for bidirectional human-AI alignment \citep{Shen2025Human}, but quantitative measurement of semantic-level accommodation across RLHF quality tiers has not been reported.

This paper addresses the following empirical question: Is RLHF alignment quality associated with systematic differences in human semantic behavior across conversations? To operationalize this question, we introduce C\_sem, a metric that measures the cosine similarity between SBERT embeddings of a human follow-up turn and its preceding AI partner turn, minus a random partner-shuffle baseline from the same tier. This subtraction is intended to isolate interaction-specific semantic alignment above chance-level topical coherence. A KNN topic-matched control (K=5) provides an additional, stricter comparison.

We compute C\_sem across the three HH-RLHF helpfulness tiers (helpful-base, helpful-rejection-sampled, helpful-online) and in both directions (human-to-AI and AI-to-human). Five sub-hypotheses are tested: (1) existence of above-baseline accommodation, (2) monotonic scaling with RLHF tier, (3) directional asymmetry favoring human-to-AI accommodation, (4) within-conversation quality discrimination as a mechanism, and (5) politeness-style mediation. Three sub-hypotheses are supported by the data; two mechanism hypotheses are falsified.
